\documentclass{article}

\usepackage{bm}


\usepackage[symbols]{glossaries}
\makeglossaries
\newglossaryentry{chemical_species}{%
  name={\ensuremath{S_i}},
  description={chemical/microbe species},
  type=symbols,
  sort=2-Si
}
\newglossaryentry{time}{%
  name={\ensuremath{t}},
  description={time},
  type=symbols,
  sort=2-t
}
\newglossaryentry{v}{%
  name={\ensuremath{v}},
  description={turnover rate},
  type=symbols,
  sort=2-v
}
\newglossaryentry{M}{%
  name={\ensuremath{M}},
  description={microbe biomass},
  type=symbols,
  sort=2-M
}
\newglossaryentry{species_index}{%
  name={\ensuremath{i}},
  description={index over chemical/microbe species},
  type=symbols,
  sort=2-i
}
\newglossaryentry{pathway_index}{%
  name={\ensuremath{j}},
  description={index over pathways},
  type=symbols,
  sort=2-j
}

\newglossaryentry{stoichiometry}{%
  name={\ensuremath{\gamma}},
  description={stoichiometric coefficient},
  type=symbols,
  sort=1-c
}
\newglossaryentry{death_rate}{%
  name={\ensuremath{\delta}},
  description={death rate},
  type=symbols,
  sort=1-d
}
\newglossaryentry{number_of_pathways}{%
  name={\ensuremath{N_P}},
  description={number of pathways},
  type=symbols,
  sort=2-NP
}
\newglossaryentry{system_size}{%
  name={\ensuremath{N_S}},
  description={system size, number of substance/microbe pools},
  type=symbols,
  sort=2-NS
}
\newglossaryentry{v_max}{%
  name={\ensuremath{v_{\mathrm{max}}}},
  description={maximum pathway turnover velocity},
  type=symbols,
  sort=2-vmax
}
\begin{document}

\glsaddallunused

\printglossary[title=List of Symbols,%
               toctitle=List of Symbols,%
               type=symbols,%
               style=list,%
               nogroupskip,%
               nonumberlist]

\newcommand{\Sdis}{S^{\mathrm{dis}}}
\newcommand{\cuej}{\mathrm{CUE}_j}%{\eta_{C,j}}

\section{Methods}

\subsection{Model}

The system is described by $N_S$ pools, representing chemical species and microbe groups changing over time.
The amount of the $i$-th substance (chemical or microbial species) that is present in the system at time $t$ is 
\begin{equation}
S_i = S_i(t).
\end{equation}

The model takes into account $N_P$ pathways, each characterized by a reaction equation and associated with a group of microbes.
The $j$-th pathway's driver are the microbes, whose current amount of biomass at time $t$ is given by
\begin{equation}
M_j(t) = S_{I^M_j}(t), 
\end{equation}
where $I^M_j$ indexes the corresponding microbe pool.

Each pathway is characterized by a reaction equation, whose stoichiometric coefficients for the $i$-th substance are given by $\gamma_{ij}$.



Microbe growth is modeled via the carbon-use-efficiency $\cuej$, indicating the fraction of carbon turnover which goes into microbe growth.

% microbes have a source of C
% => what units is biomass in? (mg/g dw or micromol/g dw)

The microbes' metabolism is composed of an anabolistic and a catabolistic fraction.
The reaction equation describes the catabolism. 
In addition, the microbes use a carbon source for the anabolism, resulting in biomass growth.
Hence, the carbon turnover of the total metabolism is composed of the stoichiometric coefficient for the respective carbon source and additional terms describing the anabolism.

The carbon use efficiency $\cuej$ indicates the anabolistic fraction with respect to the total metabolism, i.e.
\begin{equation}
\cuej = \frac{\alpha_j}{\alpha_j + \gamma_{I^M_j,j}}.
\end{equation}
Rearranging yields the effective stoichiometric contribution of the anabolism for the $j$-th pathway
\begin{equation}
\alpha_{j} = \gamma_{I^M_j,j} \frac{\cuej}{1 - \cuej}.
\end{equation}
This term describes the reduction of substance used as carbon source for the anabolism by the $j$-th pathway's microbes.
Furthermore,
\begin{equation}
\mu_{M_j,j} = \alpha_{j} N_{C,j} m_{C_j}
\end{equation}
models the resulting increase in biomass.
% TODO: \alpha_{C_j,j} is \alpha_j for correct pathway, and zero otherwise

Hence, the effective stoichiometry of the $j$-th pathway taking microbial growth into account is given by
\begin{equation}
\rho_{ij} = \gamma_{ij}  - \alpha_{C_j, j} + \mu_{M_j,j},
\end{equation}
where $\gamma_{ij}$ is the stoichiometric coefficient of the $i$-th substance in the reaction equation.







The anabolism is described via
\begin{equation}
\alpha_{M_j,j} = \gamma_{C_j,j} \frac{\cuej}{1 - \cuej}% N_{C,j} m_{C} - 
\end{equation}

The two other terms take into account microbial growth and the associated reduction of the substance used as source by the microbes.


% a = g * CUE/(1-CUE)
% a (1-CUE) = g * CUE
% a = (g + a) * CUE
% CUE = a /(g+a)
% => CUE is the fraction of anabolism in the total metabolism (g + a)


% TODO: death rate only applies to microbe pool!
The change of each substance at time $t$ is described by
\begin{equation}
\frac{\mathrm{d}S_i}{\mathrm{d}t} = \sum_{j=1}^{N_P} M_j \left(v_j \cdot \rho_{ij} - \delta_{ij}\right),
\end{equation}
taking into account contributions from the $N_{\mathrm{pwy}}$ pathways, where the time-dependency is omitted for the sake of clarity.
In order to model microbe death, the constant 
\begin{equation}
\delta_{ij} = \begin{cases} \delta & \mathrm{if}\ i = I^M_j\\
                        0 & \mathrm{otherwise},
            \end{cases}
\end{equation}
thus reducing all biomass with a constant rate.

\begin{equation}
\rho_{ij} = \gamma_{ij} + \alpha_{ij}
\end{equation}


The $j$-th pathway's turnover rate is governed by the velocity
\begin{equation}
v_j = v_{\mathrm{max},j} \cdot f_{\mathrm{MM},j} \cdot f_{\mathrm{thermo},j}\cdot f_{\mathrm{inhib},j},
\end{equation}
with the constant maximum turnover rate $v_{\mathrm{max}}$ and the time-dependent dimensionless factors $f_k$ with values between $0$ and $1$, summarizing the effects of kinetics (Michaelis-Menten and inhibition) and thermodynamics at time $t$.

The forward and inverse kinetics are modelled by the factors
\begin{equation}
f_{\mathrm{MM},j} = \prod_{i = 1}^{N_{\mathrm{S}}} \frac{S_i}{S_i + K_{m,ij}}
\end{equation}
and 
\begin{equation}
f_{\mathrm{inhib},j} = \prod_{i = 1}^{N_{\mathrm{S}}} \frac{S_i}{S_i + K_{\mathrm{inhib},ij}},
\end{equation}
where $K_{m,ij}$ and $K_{\mathrm{inhib}, ij}$ are the half-saturation coefficients with respect to the $i$-th substance in the $j$-th pathway.

If the $i$-th pool is irrelevant for the $j$-th pathway, $K_{m,ij} = 0$.
Numerics: Add $\varepsilon = 10^{-12}$ in the denominator, to avoid division by zero.


\begin{equation}
 H_{cc,i} = H_{cp}^{\circ} \cdot R\cdot T \cdot e^{C_{cp} \cdot \left(\frac{1}{T} - \frac{1}{T^{\circ}}\right)}
\end{equation}

\begin{equation}
\lambda_i = H_{cc,i} \frac{V_{\mathrm{aq}}}{V_{\mathrm{gas}}}
\end{equation}

\begin{equation}
\Sdis_i = \frac{\lambda_i}{1+\lambda_i} \cdot S_i
\end{equation}

\begin{equation}
a_i = \frac{\Sdis_i}{w}
\end{equation}

\begin{equation}
Q_j = \prod_{i =1}^{N_P} a_{i}^{\gamma_{ij}}
\end{equation}

\begin{equation}
\Delta G_{s,j} = \sum_{i = 1}^{N_P} \gamma_{ij}\Delta G_{f,i}
\end{equation}

\begin{equation}
\Delta G_{r,j} = \Delta G_{s,j} + R T \log Q_j
\end{equation}


\begin{equation}
f_{\mathrm{thermo},j} = 1 - e^{\frac{\min(0, \Delta G_{r,j} - \Delta G_{r,\min})}{RT}}
\end{equation}

$\Delta G_{r,\min} = -26 \cdot 10^3 \_ \frac{\mathrm{kJ}}{\mathrm{mol}}$


In order to avoid negative pools, 
\begin{equation}
S_i = \max(0, S_i)
\end{equation}

\end{document}
